\section{Introduction}
\label{sec:introduccion}

This work establishes a common coinductive generation of
\(\pi,\varphi,e\), the dodecaphase closure of \(\alpha\) and its
correlated analytic representation, and the algebraic composition of
the electronic coordinate. The electronic block also determines a
spin-\(1/2\) representation whose projectors give helicities
\(\pm1/2\) when a direction of propagation is specified. The same
genealogy preserves an incidence record leading to the Witt designs,
Golay codes, the Leech lattice, and \(V^\natural\). The result thus
connects determined values, recoverable memory, and algebraic
realizations through explicit operations.

Determining the values of fundamental constants constitutes a
central question in mathematical physics. A theory may establish relations
among observables, identify symmetries, and organize highly precise predictions,
while retaining parameters whose values are determined experimentally.
Understanding why those parameters take their values requires a construction
that identifies the operations, normalizations, and compatibility
conditions from which they result. This is the question organizing the present
work: the relation between the structure determining a constant, the equation
expressing that determination, and the value obtained by evaluating it.

This question is approached through a formal mathematical kernel called
\emph{Triadic Modular Holography}, hereafter HMT. Its three components are
\emph{All Possible Paths}, hereafter APP; TRIT, the name of the oriented
ternary structure; and \emph{Topological Phase Kernel}, hereafter TPK.
APP is the discrete toroidal support \(X_9=(\mathbb Z/9\mathbb Z)^2\).
Its cardinal translations define the Cayley graph of the additive
structure, with generators \(\pm(1,0),\pm(0,1)\); addition and
multiplication evaluations are computed on its vertices, retaining
remainders and quotients. TRIT is the oriented signature
\(\{-1,0,+1\}\), which distinguishes the hyperbolic, parabolic, and
elliptic regimes. TPK composes phase selection, transport of the
evaluations, and memory updating. The graph, evaluations, and dynamics
therefore have different functions on a common state.

The construction determines numerical realization together with the
structure producing it. The regions of closure, propagation, and
self-scaling are constructed in a finite domain of oriented initial
conditions; their compatible prolongation subsequently determines the coordinates
\(\pi,e,\varphi\). The dodecaphase record enters the determination of
\(\alpha\). The electronic construction composes these coordinates with a
central structure and its return operators. Finally, the incidence
retained along the route admits a realization through codes,
lattices, and a vertex algebra. Values and incidence relations are
realizations of shared data; preservation of those data under refinement
establishes their mathematical connection.

The question concerns the explanatory status of constants. When a
theory receives a parameter determined experimentally, its equations
establish the consequences that follow from that value. A structural
derivation changes this relationship: the value becomes one of the
consequences of the construction. The question of the strength of a
coupling or of a characteristic ratio of matter then becomes
the question of the mathematical relations that determine them.
The transition from parameter to consequence constitutes the guiding argument
of this article.

Numerical determination acquires, from this perspective, a constitutive
content. The equation brings together operations whose provenance must be made
explicit; its evaluation obtains the value fixed by those operations. Once
the corresponding conditions of existence, uniqueness, and compatibility
have been established, the necessity of the result belongs to the proof.
Explaining why a constant takes a value means reconstructing that
necessity within the domain of the theory. Decimal precision shows
a quantitative consequence of the structure, while the derivation
sets out the reason why precisely that consequence is obtained.

The scope of this question increases when several determinations
share a common origin. The thesis proposed by the authors links
constants to the organization that produces their reciprocal relations.
The common architecture must explain both the individuality of each
coordinate and the dependencies linking it to the others.
The order of the constructions therefore acquires demonstrative value:
it makes it possible to recognize which data are generated, which operations use them, and
which relations survive successive realizations. The diversity
of values is studied as an expression of an articulated formal unity.

In this organization, metrology performs a subsequent and
essential function. The construction fixes its results before comparing them with
observed quantities; experience examines their physical
correspondence and delimits the scope of empirical agreement. Mathematical necessity is established
with respect to the declared premises and operations, and natural
realization is examined through the identification of observables and their
independent testing. Both tasks contribute to a complete
explanation: the first determines the result within the corpus; the second
studies its relevance to describing the quantitative constitution of the
world.

This perspective makes it possible to situate the research historically. The
precedents considered below specify different ways of
relating laws, parameters, and observation. The present work adopts
as its central problem the internal determination of values and of the
structure linking them. Its arithmetic, geometric, and
algebraic developments receive their argumentative unity from that question. The
ontological aspiration lies in explaining how apparently independent
physical properties may correspond to necessary relations of a single
mathematical organization.


This problem becomes particularly clear for dimensionless couplings.
A coordinate expressed in units depends on metrological convention;
a dimensionless relation makes it possible to inquire into its determination
independently of that choice. The fine-structure constant enters
quantitatively testable relations and, at the same time, raises the question
of the structure fixing the electromagnetic coupling. The historical
program of its explanation provides a precise precedent for
studying together the value of a constant and the mathematical
organization from which it arises.

An illuminating episode appears in Dirac's 1931 work on
quantized singularities of the electromagnetic field. Its purpose
expressly addresses the existence of a minimum electric charge. The
result establishes a relation between that charge and the strength of a
magnetic monopole; Dirac notes, however, that the construction does not
by itself determine the dimensionless ratio associated with the approximate
value \(137\)\cite{dirac1931}. A genuine structural restriction is obtained,
whose demonstrative force is not equivalent to independently fixing
the coupling. This precedent shows that the question of the origin
of constants was part of fundamental research and requires
identifying exactly what each theory determines.

The history of quantum electrodynamics clarifies the difference between
that demand and the calculation of observable effects. In his Nobel lecture
of 1965, Feynman reconstructed procedures expressing level
shifts and scattering processes through the experimental mass of the
electron\cite{feynman1965}. Renormalization made it possible to obtain
finite, testable results while retaining that reference. The effectiveness
of these predictions constitutes an actual theoretical achievement; their dependence
on an experimental parameter also identifies the logical point at
which an additional determination could broaden the explanation. The
constitutive question supplements predictive capacity and studies the
conditions that fix its parameters.

The 2006 joint LEP--SLD report offers an experimental manifestation
of this architecture\cite{lep2006}. Its authors distinguish the model's
free parameters from the relations it imposes among observables.
Through radiative corrections, Z-resonance data made it possible
to infer indirectly the masses of the top quark and the W boson and to compare them
with direct measurements. Intertwined predictions and the experimental
determination of parameters formed a single testing procedure.
The relevant problem is therefore to identify the dependencies that a
construction determines internally and those whose specification continues
to require empirical information.

In this context, \emph{epistemological determination} will mean
the procedure that establishes a value through the relation among theory,
observation, and inference; \emph{constitutive determination} will mean the
deduction of that value from the relations defining the object
under consideration. The latter expression here designates a mathematical requirement:
to make explicit the domain, operators, normalizations, and
compatibilities that select the result. Its physical relevance
arises from subsequent comparison with observed quantities.
The change in perspective affects the order of explanation: the value is
examined as a consequence of a determined structure. The ontological
aspiration consists in relating that structure to the constitution of the
physical object; mathematical proof and experimental testing establish
the scope within which such a relation may be sustained.


The fine-structure constant makes the scope of the
problem particularly visible. In the usual electromagnetic description, \(\alpha\) is a
dimensionless coupling: it expresses a physical relation that remains when
units change. Pauli linked the explanation of its value to
understanding the atomistic structure of electricity\cite{pauli1946}.
Here the mathematical construction of the coupling is presented first,
followed by that of the electron to which this interaction is attributed. Their proximity
in the argument responds to a structural dependency: the data
determining the coupling also enter the composition of the
electronic coordinate. Physical comparison finally examines these
previously fixed expressions.

The historical perspective also encompasses the relation among action, phase,
and orientation. In his path formulation, Feynman assigns trajectories
a phase dependent on action and obtains the amplitude by
composing their contributions\cite{feynman1948}. The name
\emph{All Possible Paths} takes up this conceptual reference. In the present
construction, paths are defined on a finite arithmetic structure;
the admissibility rules and their evaluations are specified through
integer operations. The relation to physical action is formulated at the
dimensional stage, once that domain has been constructed.

The link between action and angular phase also appears in
Dirac's Lagrangian treatment\cite{dirac1933}. The expression \(\hbar=h/(2\pi)\)
places action relative to an angular coordinate in radians. The present development
first studies the determinants, quotients, and transports that
produce a dimensionless action scale. Its application to a physical
unit is distinguished from the ratios in which that common scale cancels.
This separation is decisive in explaining electronic composition
and the role of chronological return: the clock quantum provides a
scale, and the electronic section preserves its structural identity.

In his 1965 Nobel lecture, Feynman recalled a conversation with
John Archibald Wheeler about the identity of the charge and mass of
electrons\cite{feynman1965}. Wheeler imagined a single worldline
which, traversing time with alternating orientations, would appear in a
temporal section as multiple electrons and positrons. Feynman distinguished
this one-electron conjecture from the representation he incorporated into
his work: a positron as a segment of an electronic line with reversed temporal
orientation. His 1949 article develops this representation and
acknowledges Stückelberg's precedent\cite{feynman1949}. This episode
historically situates the problem of trajectory orientation. In HMT,
bidirectionality is expressed through involutions, transports, and
return conditions whose domain is specified in each construction.

The arithmetic content begins with the coexistence of two evaluations.
Each row of the multiplication table is obtained by additive accumulation;
modular reduction records its residues, while quotients preserve
the multiples of the base crossing the boundary. The difference between
addition and multiplication thus acquires an exact expression on a common support.
Mixed compositions may be order-sensitive and give rise to
discrete curvature. TRIT distinguishes the elliptic, parabolic, and
hyperbolic regimes through their quadratic relations, and the exponential provides
a common law for all three. This organization makes it possible to study a number
as the scalar realization of a compatible section preserving more data
than its isolated real coordinate.

Refinement makes this distinction explicit. A compatible sequence of
cylinders of decreasing diameter determines an Archimedean coordinate;
the trajectory simultaneously preserves orientation, transport quotients,
regime, and memory. Nonadic holonomy identifies phase return,
while its lift records the advance in depth. The Sturmian
representation describes the aperiodic order of this prolongation and provides
an operator formulation of its two phases. The scope of this symbolic
construction is established by its operators and invariants.

The record \(K\) receives its own definition before entering
\(\alpha\). It is constructed by a reversible integral transformation of the
signed record; its rational realization \(\kappa\) preserves the twelve
components, origin, and orientation fixed. Cyclic differences
determine the sectoral relations, and the uniform component completes
its reconstruction. For \(\alpha\), positional determination
by integer normalization is distinguished from analytic determination through the vacancy
resolvent. Their identification requires preserving compatibility of the
prolongations. The consequences for \(e+\pi\) will be expressed in terms
of these exact relations and of the arithmetic conditions that actually
allow decisions to be made.

The unique selection of \(K\) is established through operations on
explicit populations. The sequence
\(40320\to21902\to19446\to79\to1\) passes from orderings to
indexed grid incidences, distinct numerators, candidates on the
exceptional face, and heterotypic survivors. The regional count
\(196\cdot197\cdot196\to331\to2\to1\) applies Gram, occupancy,
distance, and orientation conditions to blocks from three catalogues.
Section~\ref{act21:sec:relacion-censos} distinguishes the domains of
the two counts; the selection and inversion proofs specify
how their outputs compose with the reversible record.

The electronic section develops the origin of each factor in the equation:
the matrix norm \(\sqrt3/4\), the regional transfer
\(\varphi/\pi^2\), the cubic term in \(\alpha\), torsional
normalization, and the multiplicative memory of return. Two readers of the
central trajectory determine the same coefficient triple through
different operations. The same recurrence determines a \(27\)-cell support
and two oriented histories distinguished by the integral chart. Its harmonic
reading identifies three modes of the central record and expresses the
return coefficients through their spectral powers. The electronic equation
thus retains an explicit dependence on the generated trajectory.
The local structure and the complete coordinate are
retained during dimensional realization. Comparison is presented
with the CODATA 2018 and 2022 adjustments\cite{codata2018,codata2022}, showing
central values, their uncertainties, and the calculated differences.

The exceptional prolongation takes up the boundary data of that same
route. Changes of representation pass from ternary sequences to
codes, from supports to incidence designs, and from discriminant classes to
lattices. Construction of the Leech neighbor makes it possible to apply the
Frenkel, Lepowsky, and Meurman procedure to obtain \(V^\natural\);
the Monster acts as the automorphism group of that algebra, and Moonshine
describes its graded traces\cite{flm,borcherds}. Scalar realization
and incidence realization constitute the two projections examined
in the article: both start from identified data and preserve the
pertinent relations of the construction.

% BEGIN S27_FRAMING
Two developments specify the articulation between integer records
and their incidences. Section~\ref{carry27:seccion} derives the carry
cocycle from Euclidean division of the regional columns. The complete
table of \(27\) inputs, polynomial representations, and rank proofs
establish the algebraic contribution of the quotients. The transport
identities retain this contribution together with the residue;
the cylinder bounds determine compatibility of the decimal reading.

Section~\ref{star27:comparacion} then compares the stars of
\(B_K=H_e\cap P_\pi\), \(B_\varphi=H_\varphi\cap P_\pi\), and
\(B_{\rm APP}=H_{\rm APP}\cap P_\pi\), polarized by the negative
support of the pre-carry cochain. Their twelve completions, the census
of the \(495\) four-subsets, and the stabilizers of the ordered frame
specify the information retained by this incidence reading.
The connection between the two developments lies in the regional
and oriented records produced by APP--TRIT--TPK: integer evaluation
retains magnitudes and carries, and boundary realization retains
supports, intersections, and symmetries.

% END S27_FRAMING


% BEGIN R27_FRAMING_INTRO_EN
The finite construction specifies which information is transmitted at
each stage. The alphabet of one window consists of \(42\) blocks obtained
from the emitter's two signatures; the six windows produce the catalogue
on which regional junctions are realized. Two records with the same
histogram and different temporal correlation identify a contribution of
order omitted by the marginal reading. The four gauges of the
multiplicative observable and the closure controls are examined on fixed
domains, preserving the provenance of the register and its boundary
conditions.

The same discipline distinguishes three extensions of APP: changing the
modulus, periodically repeating the support, and increasing the number
of coordinates. Multiplicative compression into three classes preserves
an index-two sublattice whose rank-two component realizes a Pell
iteration. Encoding pairs of trits transports addition through its carry.
Subsequently, quadratic correlation constructs the identity of the
six-position chart; dual neutrality reduces the eight saturated boundaries
to two, and the fixed orientation selects their representative. The
level-three divisibility result belongs at the end of the exceptional
realization, where its proof uses the rational identity already established.
% END R27_FRAMING_INTRO_EN

% BEGIN R28_FRAMING_INTRO_EN
The regional linkage is further developed on the emitter's labelled
multigraph and its phase quotient. Five walks of minimum length eight
join the APP, propagation, and autoscale anchors in each closure region.
Their oriented witnesses and finite minimality proof complete the
account of the junctions and of the information retained by each
representation.
% END R28_FRAMING_INTRO_EN


The enumeration of initial conditions specifies the admissible states;
the computational certificates allow their finite relations to be checked,
and the provenance records identify the data for each operation.
The definitions, lemmas, and their proofs establish continuity between
these objects. The passages to the limit retain their hypotheses and
restriction maps.
The construction in Section~\ref{subsec:alpha-publicaciones-completas}
specifies every coefficient of the correlated analytic representation and
proves its compatibility with the integer closure. Equality of the prefixes
of both representations at every depth is established in
Theorem~\ref{thm:alpha-dos-publicaciones-completas}, within the geometric
class and domain defined there.
The guiding question remains throughout
the development: which structure determines the value, and which relations persist
when its representation changes.

